% Explanation certificate for heart-s2-c2 (genome d4bbc960-816a-4781-8979-fdd4165debf2).
% Generated by explaneat export (schema 1.4.0) from the JSON record alone;
% content sha256 1276e9a13f62dbc089618dbbcbf8c7fcbf1ae0d7094c9016ac0b90a0a0c036ac.
% Drafter recorded in the JSON: ai_agent (claude). Do not edit by hand: re-render with
%   uv run python -m explaneat export render --from heart-s2-c2.json
% TikZ modules are input from \explaneatcertdir{<label>} (empty by default: same
% directory; index.tex sets it to the genome's subdirectory).
\providecommand{\explaneatcertdir}[1]{}
\begin{drafted}[appendix: heart-s2-c2 explanation certificate]
\section{Explanation certificate: heart-s2-c2}\label{app:heart-s2-c2:sec}
This section is the explanation certificate of the heart-s2-c2 network. It is rendered from the canonical JSON record \texttt{heart-s2-c2.json} (schema 1.4.0, SHA-256 \texttt{1276e9a13f62dbc0}\ldots), which also holds every evidence payload, the full graphs with weights, and the explanation record. \Cref{app:heart-s2-c2:summary} summarises the network and its coverage.

\begin{table}[htbp]
\centering\small
\begin{tabularx}{\linewidth}{@{}lX@{}}
\toprule
Dataset & Heart Disease (binary, prepared) \\
Experiment & 10seed-heart-s2, seed 2 \\
Genome & \texttt{d4bbc960-816a-4781-8979-fdd4165debf2} (NEAT key 9685, generation 39) \\
Explanation & \texttt{3531d96c-aa00-4141-863a-2c85155afe8f} (31 operations) \\
Train AUC / accuracy & 0.9403 / 0.8967 (n = 242) \\
Test AUC / accuracy & 0.9009 / 0.8197 (n = 61) \\
Output activation & sigmoid \\
Phenotype & 31 nodes (26 inputs, 4 hidden, 1 output), 24 enabled edges \\
Explained model state & 33 nodes, 24 enabled edges \\
Refactoring operations & split\_\allowbreak{}node 2 \\
Function preservation & max $|\Delta\,\mathrm{output}|$ = $2.22\times 10^{-16}$ on the training split \\
Annotations & 3 (2 leaf, 1 composition) \\
Structural coverage & 1 (32 of 32 candidate nodes; 8 Null) \\
Compositional coverage & 1 \\
Uncovered nodes / edges & 0 / 0 \\
Missing compositions & 0 \\
Full coverage & yes \\
Drafted by & AI-assisted draft \\
Verified by & unverified \\
Record hash (SHA-256) & \texttt{1276e9a13f62dbc0}\ldots \\
\bottomrule
\end{tabularx}
\caption{Explanation certificate summary: heart-s2-c2.}
\label{app:heart-s2-c2:summary}
\end{table}

\paragraph{Drafter's summary.} Mostly logistic regression: 16 direct edges (logit SD 3.66; train AUC 0.934 alone, L2 logistic regression 0.9355) plus one consequential nonlinear part to 0.940. The protective gate 3842 = ReLU(1.02 ReLU(0.688 thalach - 0.374 thal:6 + 2.15 [non-anginal pain] - 0.347) - 0.295), read out -0.519, opens for nearly all non-anginal-pain patients, for others only above about +0.8 SD max heart rate, and never with a fixed thallium defect; once open a higher max heart rate lowers risk 3.4x as steeply. A restecg:2 relay is negligible. Unlike the generation-38 heart-s2 record (age/thalach/thal:7 gate, flat-slope gate, oldpeak hinges), age and oldpeak here act only linearly and the gate is keyed on chest-pain type and fixed defect instead.

\paragraph{What resisted explanation.} Nothing resisted coverage; only input fan-out needed splitting (thalach, thal:6). Relay 3607 is a step on the binary restecg:2 input, so exactly affine on the data and negligible (read out -0.021). Eight inputs are unconnected (thal:3, thal:4, slope:3, exang:0, exang:1, restecg:0, restecg:1, fbs:0) and listed as zero-effect root entries; exercise angina is ignored entirely. Counter-directional signs reported, not explained: cp:2 +0.26 and fbs:1 -0.22 in the linear core, and thal:6 (fixed defect) -0.30 directly plus closing the protective gate.

\subsection{Plain-language summary}\label{app:heart-s2-c2:plain}
\noindent{\small\textit{AI-assisted draft for review (unverified). Each sentence cites the claims (C01, C02, \ldots) that back it; every cited claim was computed on the model and the training split and passed its automatic checks. Magnitudes are changes in predicted risk for a typical patient; confidence words follow the test-split check and n. The text describes what the model does, not the world.}}\par\smallskip
\paragraph{Summary (drafted).} For a typical patient (a 56-year-old man with no chest pain symptoms and median test results) the model predicts a risk of heart disease of about 32\%; it adds up most of its inputs independently, and two parts add to that sum: one acts on 104 of 242 training patients through chest pain type, maximum heart rate and the thallium scan, and one through the resting ECG barely changes the predictions.\,{\scriptsize\textcolor{NavyBlue}{[C29, C37, C41]}} These add up independently, and these lower the risk: through this sum alone, compared with no chest pain symptoms, typical angina takes a typical patient's predicted risk from about 32\% to about 2\%, non-anginal chest pain to about 7\% and atypical angina to about 14\%; being a woman lowers it to about 4\% and fasting blood sugar above 120 mg/dL to about 20\%; and a higher maximum heart rate lowers it steadily, from about 37\% at 118 to about 28\% at 175 beats per minute (chest pain type and maximum heart rate also act through the heart rate gate, below).\,{\scriptsize\textcolor{NavyBlue}{[C04, C12, C16, C08, C18, C20]}} Each extra vessel coloured by fluoroscopy adds about the same, from about 32\% at 0 to about 96\% at 2 vessels, flattening out near the top because the predicted risk cannot go much higher; a flat rather than upsloping ST segment raises it to about 81\%; and ST depression on exercise, cholesterol and resting blood pressure each raise it steadily, from about 25\% at 0 mm to about 55\% at 2.8 mm, from about 22\% at 193 mg/dL to about 50\% at 308 mg/dL, and from about 21\% at 110 mm Hg to about 48\% at 155 mm Hg.\,{\scriptsize\textcolor{NavyBlue}{[C01, C02, C03, C07, C09]}} Through this sum alone a reversible defect on the thallium scan raises it to about 57\% (the thallium scan also acts through the heart rate gate, below); older age raises it steadily but only from about 31\% at 42 to about 32\% at 66 years; and the model ignores chest pain brought on by exercise.\,{\scriptsize\textcolor{NavyBlue}{[C10, C21, C28]}} Some rarer results also count (the model has little data on this, n = 15 to 17): a downsloping ST segment raises the risk to about 54\%, and through this sum alone a fixed defect on the thallium scan lowers it to about 12\%.\,{\scriptsize\textcolor{NavyBlue}{[C13, C14]}} This part reads chest pain type, maximum heart rate and the thallium scan, and acts on 104 of 242 training patients, 24 of whom had heart disease; for the typical patient it acts on, the predicted risk is about 5\% without it and about 2\% with it.\,{\scriptsize\textcolor{NavyBlue}{[C37]}} Across the whole model, non-anginal chest pain lowers the predicted risk, and even more so without a fixed defect on the thallium scan: from about 32\% to about 3\% then, against about 12\% to about 2\% with one (the model has little data on this, n = 15).\,{\scriptsize\textcolor{NavyBlue}{[C33]}} Across the whole model, a higher maximum heart rate lowers the risk, and even more so with non-anginal chest pain, where it falls steadily from about 7\% at 118 to about 2\% at 175 beats per minute, flattening out near the bottom because the predicted risk cannot go much lower; for other chest pain types it falls from about 37\% to about 26\%.\,{\scriptsize\textcolor{NavyBlue}{[C35]}} Through this part alone, for most patients a higher maximum heart rate does nothing until about 168 beats per minute, then lowers the risk steadily, but for patients with non-anginal chest pain it lowers it steadily all along: for a typical patient from about 32\% at 118 to about 29\% at 175 beats per minute (the model has little data on this, n = 14).\,{\scriptsize\textcolor{NavyBlue}{[C31, C36]}} This part reads only the resting ECG: signs of left ventricular hypertrophy take a typical patient's predicted risk from about 32\% to about 31\%, and switching the part off changes any prediction by at most about half a percentage point.\,{\scriptsize\textcolor{NavyBlue}{[C39, C41]}}
\paragraph{The whole model.} The model adds up its inputs into one score: the vessels coloured (each extra vessel adding about the same), the ST segment, chest pain type and sex each move the predicted risk on their own, and two parts add to that sum: a heart rate gate that lowers the risk, and a small step for the resting ECG.\,{\scriptsize\textcolor{NavyBlue}{[C29, C01, C02, C05, C08]}} \textit{Evidence.} It is made up of the inputs that feed the score directly and two parts, the heart rate gate and the resting ECG step (each explained below); the score is their sum, which was checked on all 242 training patients.\,{\scriptsize\textcolor{NavyBlue}{[C29]}} Chest pain brought on by exercise feeds nothing, so the model ignores it.\,{\scriptsize\textcolor{NavyBlue}{[C28]}} Across the whole model, compared with no chest pain symptoms, typical angina takes a typical patient's predicted risk from about 32\% to about 2\%, non-anginal chest pain to about 3\% and atypical angina to about 14\%.\,{\scriptsize\textcolor{NavyBlue}{[C05, C06, C17]}} Across the whole model, a higher maximum heart rate lowers the risk, for most patients slowly at first, then faster, but for patients with non-anginal chest pain steadily: from about 37\% at 118 to about 26\% at 175 beats per minute, against about 28\% through the sum alone.\,{\scriptsize\textcolor{NavyBlue}{[C19, C20]}} Across the whole model, a reversible defect on the thallium scan raises it to about 57\%, and a fixed defect lowers it to about 12\% (the model has little data on this, n = 15).\,{\scriptsize\textcolor{NavyBlue}{[C11, C15]}}
\paragraph{Heart rate gate.} This part lowers the predicted risk for patients with non-anginal chest pain or a high maximum heart rate, unless the thallium scan shows a fixed defect (the model has little data on this, n = 15).\,{\scriptsize\textcolor{NavyBlue}{[C37, C33]}} \textit{Evidence.} It acts on 104 of 242 training patients, 24 of whom had heart disease; for the typical patient it acts on, the predicted risk is about 5\% without it and about 2\% with it.\,{\scriptsize\textcolor{NavyBlue}{[C37]}} Through it, non-anginal chest pain takes a typical patient's predicted risk from about 32\% to about 16\%, and for the typical patient it acts on a fixed defect on the thallium scan raises it from about 2\% to about 4\% (the model has little data on this, n = 15).\,{\scriptsize\textcolor{NavyBlue}{[C30, C32]}} Through it, for most patients a higher maximum heart rate does nothing until about 168 beats per minute, then lowers the risk steadily, but for patients with non-anginal chest pain it lowers it steadily all along: for a typical patient from about 32\% at 118 to about 29\% at 175 beats per minute (the model has little data on this, n = 14).\,{\scriptsize\textcolor{NavyBlue}{[C31, C36]}} Across the whole model, a higher maximum heart rate lowers the risk more without a fixed defect on the thallium scan, for most patients slowly at first, then faster, but for patients with non-anginal chest pain steadily: from about 37\% to about 26\%, against about 14\% to about 10\% with one (the model has little data on this, n = 15).\,{\scriptsize\textcolor{NavyBlue}{[C34]}} Its formula reproduces the part exactly on all 242 training patients.\,{\scriptsize\textcolor{NavyBlue}{[C38]}}
\paragraph{Resting ECG step.} This part barely changes the predictions: signs of left ventricular hypertrophy on the resting ECG take a typical patient's predicted risk from about 32\% to about 31\%.\,{\scriptsize\textcolor{NavyBlue}{[C39, C41]}} \textit{Evidence.} Switching it off changes any training patient's predicted risk by at most about half a percentage point.\,{\scriptsize\textcolor{NavyBlue}{[C41]}} Its formula reproduces the part exactly on all 242 training patients.\,{\scriptsize\textcolor{NavyBlue}{[C42]}}
\paragraph{Worth checking against what is known.} Worth checking: going from 0 to 2 vessels coloured moves a typical patient's predicted risk from about 32\% to about 96\%.\,{\scriptsize\textcolor{NavyBlue}{[C01]}} Worth checking: patients with typical angina get the lowest predicted risk of all chest pain types, below patients with no chest pain symptoms.\,{\scriptsize\textcolor{NavyBlue}{[C05, C06, C17]}} Worth checking: the model treats fasting blood sugar above 120 mg/dL as lowering the risk of heart disease (for a typical patient from about 32\% to about 20\%).\,{\scriptsize\textcolor{NavyBlue}{[C18]}} Worth checking: the model treats a fixed defect on the thallium scan as lowering the risk (for a typical patient from about 32\% to about 12\%), from only 15 training patients.\,{\scriptsize\textcolor{NavyBlue}{[C15]}} Worth checking: the model ignores whether exercise brings on chest pain.\,{\scriptsize\textcolor{NavyBlue}{[C28]}}
\paragraph{Claims.}
\par{\tiny\noindent
\begin{tabularx}{\linewidth}{@{}l>{\raggedright\arraybackslash}X>{\raggedright\arraybackslash}p{0.11\linewidth}>{\raggedright\arraybackslash}p{0.16\linewidth}>{\raggedright\arraybackslash}p{0.05\linewidth}>{\raggedright\arraybackslash}p{0.15\linewidth}l>{\raggedright\arraybackslash}p{0.08\linewidth}@{}}
\toprule
Claim & What is checked & Risk, typical patient & Estimate (log-odds; OR); range across patients & n & Test check & Verified & Words \\
\midrule
C01 & direction: number of vessels coloured by fluoroscopy: 0 vessels to 2 vessels (whole model) \newline\textit{shape: steady (equal steps); 100\% of patients same; verified; risk saturates near 100\%} & 32\% to 96\% & log-odds 3.99, OR 54.0; range across patients 3.99 to 3.99 & 242 & train; test unclear (AUC gain 0.037 [-0.038, 0.119], n=61) & \textcolor{OliveGreen}{yes} & on the training data only \\
C02 & direction: slope of the ST segment at peak exercise: an upsloping ST segment to a flat ST segment (whole model) \newline\textit{shape: categorical; highest to lowest: a flat ST segment, a downsloping ST segment, an upsloping ST segment} & 32\% to 81\% & log-odds 2.21, OR 9.14; range across patients 2.21 to 2.21 & 115, 110 & train; test disagrees (AUC gain -0.004 [-0.044, 0.034], n=61) & \textcolor{OliveGreen}{yes} & on the training data only \\
C03 & direction: ST depression on exercise: 0.0 mm to 2.8 mm (whole model) \newline\textit{shape: steady; 100\% of patients same; verified} & 25\% to 55\% & log-odds 1.31, OR 3.72; range across patients 1.31 to 1.31 & 242 & train; test unclear (AUC gain 0.011 [-0.005, 0.034], n=61) & \textcolor{OliveGreen}{yes} & on the training data only \\
C04 & direction: chest pain type: no chest pain (asymptomatic) to typical angina (direct contribution only, not through any part) \newline\textit{shape: categorical; highest to lowest: no chest pain (asymptomatic), atypical angina, non-anginal chest pain, typical angina} & 32\% to 2\% & log-odds -3.03, OR 0.05; range across patients -3.03 to -3.03 & 108, 22 & train; test agrees (AUC gain 0.039 [0.009, 0.076], n=61) & \textcolor{OliveGreen}{yes} & clearly \\
C05 & direction: chest pain type: no chest pain (asymptomatic) to typical angina (whole model) \newline\textit{shape: categorical; highest to lowest: no chest pain (asymptomatic), atypical angina, non-anginal chest pain, typical angina} & 32\% to 2\% & log-odds -3.03, OR 0.05; range across patients -3.03 to -3.03 & 108, 22 & train; test agrees (AUC gain 0.057 [0.017, 0.104], n=61) & \textcolor{OliveGreen}{yes} & clearly \\
C06 & direction: chest pain type: no chest pain (asymptomatic) to non-anginal chest pain (whole model) \newline\textit{shape: categorical; highest to lowest: no chest pain (asymptomatic), atypical angina, non-anginal chest pain, typical angina} & 32\% to 3\% & log-odds -2.53, OR 0.08; range across patients -2.88 to -1.96 & 108, 70 & train; test agrees (AUC gain 0.057 [0.017, 0.104], n=61) & \textcolor{OliveGreen}{yes} & clearly \\
C07 & direction: serum cholesterol: 193 mg/dL (5.0 mmol/L) to 308 mg/dL (8.0 mmol/L) (whole model) \newline\textit{shape: steady; 100\% of patients same; verified} & 22\% to 50\% & log-odds 1.28, OR 3.58; range across patients 1.28 to 1.28 & 242 & train; test disagrees (AUC gain -0.011 [-0.039, 0.012], n=61) & \textcolor{OliveGreen}{yes} & on the training data only \\
C08 & direction: sex: male to female (whole model) \newline\textit{shape: categorical; highest to lowest: male, female} & 32\% to 4\% & log-odds -2.34, OR 0.10; range across patients -2.34 to -2.34 & 159, 83 & train; test unclear (AUC gain 0.027 [-0.01, 0.074], n=61) & \textcolor{OliveGreen}{yes} & on the training data only \\
C09 & direction: resting blood pressure: 110 mm Hg to 155 mm Hg (whole model) \newline\textit{shape: steady; 100\% of patients same; verified} & 21\% to 48\% & log-odds 1.22, OR 3.38; range across patients 1.22 to 1.22 & 242 & train; test unclear (AUC gain 0.009 [-0.013, 0.034], n=61) & \textcolor{OliveGreen}{yes} & on the training data only \\
C10 & direction: thallium stress test result: thallium stress test result other than a fixed defect on the thallium scan, a reversible defect on the thallium scan to a reversible defect on the thallium scan (direct contribution only, not through any part) \newline\textit{shape: categorical; highest to lowest: a reversible defect on the thallium scan, thallium stress test result other than a fixed defect on the thallium scan, a reversible defect on the thallium scan, a fixed defect on the thallium scan} & 32\% to 57\% & log-odds 1.04, OR 2.82; range across patients 1.04 to 1.04 & 137, 90 & train; test unclear (AUC gain 0.018 [-0.012, 0.053], n=61) & \textcolor{OliveGreen}{yes} & on the training data only \\
C11 & direction: thallium stress test result: thallium stress test result other than a fixed defect on the thallium scan, a reversible defect on the thallium scan to a reversible defect on the thallium scan (whole model) \newline\textit{shape: categorical; highest to lowest: a reversible defect on the thallium scan, thallium stress test result other than a fixed defect on the thallium scan, a reversible defect on the thallium scan, a fixed defect on the thallium scan} & 32\% to 57\% & log-odds 1.04, OR 2.82; range across patients 1.04 to 1.04 & 137, 90 & train; test unclear (AUC gain 0.018 [-0.012, 0.053], n=61) & \textcolor{OliveGreen}{yes} & on the training data only \\
C12 & direction: chest pain type: no chest pain (asymptomatic) to non-anginal chest pain (direct contribution only, not through any part) \newline\textit{shape: categorical; highest to lowest: no chest pain (asymptomatic), atypical angina, non-anginal chest pain, typical angina} & 32\% to 7\% & log-odds -1.75, OR 0.17; range across patients -1.75 to -1.75 & 108, 70 & train; test agrees (AUC gain 0.039 [0.009, 0.076], n=61) & \textcolor{OliveGreen}{yes} & clearly \\
\bottomrule
\end{tabularx}}
\par{\tiny\noindent
\begin{tabularx}{\linewidth}{@{}l>{\raggedright\arraybackslash}X>{\raggedright\arraybackslash}p{0.11\linewidth}>{\raggedright\arraybackslash}p{0.16\linewidth}>{\raggedright\arraybackslash}p{0.05\linewidth}>{\raggedright\arraybackslash}p{0.15\linewidth}l>{\raggedright\arraybackslash}p{0.08\linewidth}@{}}
\toprule
Claim & What is checked & Risk, typical patient & Estimate (log-odds; OR); range across patients & n & Test check & Verified & Words \\
\midrule
C13 & direction: slope of the ST segment at peak exercise: an upsloping ST segment to a downsloping ST segment (whole model) \newline\textit{shape: categorical; highest to lowest: a flat ST segment, a downsloping ST segment, an upsloping ST segment} & 32\% to 54\% & log-odds 0.94, OR 2.56; range across patients 0.94 to 0.94 & 115, 17 (small n) & train; test disagrees (AUC gain -0.004 [-0.044, 0.034], n=61) & \textcolor{OliveGreen}{yes} & the model has little data on this; on the training data only \\
C14 & direction: thallium stress test result: thallium stress test result other than a fixed defect on the thallium scan, a reversible defect on the thallium scan to a fixed defect on the thallium scan (direct contribution only, not through any part) \newline\textit{shape: categorical; highest to lowest: a reversible defect on the thallium scan, thallium stress test result other than a fixed defect on the thallium scan, a reversible defect on the thallium scan, a fixed defect on the thallium scan} & 32\% to 12\% & log-odds -1.24, OR 0.29; range across patients -1.24 to -1.24 & 137, 15 (small n) & train; test unclear (AUC gain 0.018 [-0.012, 0.053], n=61) & \textcolor{OliveGreen}{yes} & the model has little data on this; on the training data only \\
C15 & direction: thallium stress test result: thallium stress test result other than a fixed defect on the thallium scan, a reversible defect on the thallium scan to a fixed defect on the thallium scan (whole model) \newline\textit{shape: categorical; highest to lowest: a reversible defect on the thallium scan, thallium stress test result other than a fixed defect on the thallium scan, a reversible defect on the thallium scan, a fixed defect on the thallium scan} & 32\% to 12\% & log-odds -1.00, OR 0.37; range across patients -1.24 to -0.42 & 137, 15 (small n) & train; test unclear (AUC gain 0.018 [-0.012, 0.053], n=61) & \textcolor{OliveGreen}{yes} & the model has little data on this; on the training data only \\
C16 & direction: chest pain type: no chest pain (asymptomatic) to atypical angina (direct contribution only, not through any part) \newline\textit{shape: categorical; highest to lowest: no chest pain (asymptomatic), atypical angina, non-anginal chest pain, typical angina} & 32\% to 14\% & log-odds -1.05, OR 0.35; range across patients -1.05 to -1.05 & 108, 42 & train; test agrees (AUC gain 0.039 [0.009, 0.076], n=61) & \textcolor{OliveGreen}{yes} & clearly \\
C17 & direction: chest pain type: no chest pain (asymptomatic) to atypical angina (whole model) \newline\textit{shape: categorical; highest to lowest: no chest pain (asymptomatic), atypical angina, non-anginal chest pain, typical angina} & 32\% to 14\% & log-odds -1.05, OR 0.35; range across patients -1.05 to -1.05 & 108, 42 & train; test agrees (AUC gain 0.057 [0.017, 0.104], n=61) & \textcolor{OliveGreen}{yes} & clearly \\
C18 & direction: fasting blood sugar: fasting blood sugar of 120 mg/dL or less to fasting blood sugar above 120 mg/dL (whole model) \newline\textit{shape: categorical; highest to lowest: fasting blood sugar of 120 mg/dL or less, fasting blood sugar above 120 mg/dL} & 32\% to 20\% & log-odds -0.65, OR 0.52; range across patients -0.65 to -0.65 & 209, 33 & train; test unclear (AUC gain 0.000 [-0.011, 0.013], n=61) & \textcolor{OliveGreen}{yes} & on the training data only \\
C19 & direction: maximum heart rate: 118 beats per minute to 175 beats per minute (whole model) \newline\textit{shape: accelerating; 67\% of patients same; mixed (but steady for patients with non-anginal chest pain)} & 37\% to 26\% & log-odds -0.73, OR 0.48; range across patients -1.32 to -0.51 & 242 & train; test agrees (AUC gain 0.014 [0.002, 0.037], n=61) & \textcolor{OliveGreen}{yes} & clearly \\
C20 & direction: maximum heart rate: 118 beats per minute to 175 beats per minute (direct contribution only, not through any part) \newline\textit{shape: steady; 100\% of patients same; verified} & 37\% to 28\% & log-odds -0.39, OR 0.68; range across patients -0.39 to -0.39 & 242 & train; test unclear (AUC gain 0.006 [0.000, 0.020], n=61) & \textcolor{OliveGreen}{yes} & on the training data only \\
C21 & direction: age: 42 years to 66 years (whole model) \newline\textit{shape: steady; 100\% of patients same; verified} & 31\% to 32\% & log-odds 0.07, OR 1.08; range across patients 0.07 to 0.07 & 242 & train; test disagrees (AUC gain -0.001 [-0.005, 0.000], n=61) & \textcolor{OliveGreen}{yes} & slightly; on the training data only \\
C22 & direction: resting ECG: resting ECG other than signs of left ventricular hypertrophy on the resting ECG to signs of left ventricular hypertrophy on the resting ECG (whole model) \newline\textit{shape: categorical; highest to lowest: resting ECG other than signs of left ventricular hypertrophy on the resting ECG, signs of left ventricular hypertrophy on the resting ECG} & 32\% to 31\% & log-odds -0.02, OR 0.98; range across patients -0.02 to -0.02 & 123, 119 & train; test unclear (AUC gain 0.000 [0.000, 0.000], n=61) & \textcolor{OliveGreen}{yes} & slightly; on the training data only \\
C23 & magnitude: number of vessels coloured by fluoroscopy: 0 vessels to 2 vessels, typical patient \newline\textit{shape: steady (equal steps); 100\% of patients same; verified; risk saturates near 100\%} & 32\% to 96\% & log-odds 3.99, OR 54.0 & 242 & train; test unclear (AUC gain 0.037 [-0.038, 0.119], n=61) & \textcolor{OliveGreen}{yes} & on the training data only \\
C24 & magnitude: slope of the ST segment at peak exercise: an upsloping ST segment to a flat ST segment, typical patient \newline\textit{shape: categorical; highest to lowest: a flat ST segment, a downsloping ST segment, an upsloping ST segment} & 32\% to 81\% & log-odds 2.21, OR 9.14 & 115, 110 & train; test disagrees (AUC gain -0.004 [-0.044, 0.034], n=61) & \textcolor{OliveGreen}{yes} & on the training data only \\
\bottomrule
\end{tabularx}}
\par{\tiny\noindent
\begin{tabularx}{\linewidth}{@{}l>{\raggedright\arraybackslash}X>{\raggedright\arraybackslash}p{0.11\linewidth}>{\raggedright\arraybackslash}p{0.16\linewidth}>{\raggedright\arraybackslash}p{0.05\linewidth}>{\raggedright\arraybackslash}p{0.15\linewidth}l>{\raggedright\arraybackslash}p{0.08\linewidth}@{}}
\toprule
Claim & What is checked & Risk, typical patient & Estimate (log-odds; OR); range across patients & n & Test check & Verified & Words \\
\midrule
C25 & magnitude: ST depression on exercise: 0.0 mm to 2.8 mm, typical patient \newline\textit{shape: steady; 100\% of patients same; verified} & 25\% to 55\% & log-odds 1.31, OR 3.72 & 242 & train; test unclear (AUC gain 0.011 [-0.005, 0.034], n=61) & \textcolor{OliveGreen}{yes} & on the training data only \\
C26 & magnitude: chest pain type: no chest pain (asymptomatic) to typical angina, typical patient \newline\textit{shape: categorical; highest to lowest: no chest pain (asymptomatic), atypical angina, non-anginal chest pain, typical angina} & 32\% to 2\% & log-odds -3.03, OR 0.05 & 108, 22 & train; test agrees (AUC gain 0.057 [0.017, 0.104], n=61) & \textcolor{OliveGreen}{yes} & clearly \\
C27 & magnitude: serum cholesterol: 193 mg/dL (5.0 mmol/L) to 308 mg/dL (8.0 mmol/L), typical patient \newline\textit{shape: steady; 100\% of patients same; verified} & 22\% to 50\% & log-odds 1.28, OR 3.58 & 242 & train; test disagrees (AUC gain -0.011 [-0.039, 0.012], n=61) & \textcolor{OliveGreen}{yes} & on the training data only \\
C28 & null: model ignores exercise-induced angina (exact rule) &  & max change 0.000 (max $|$change in log-odds$|$) & 242 & train & \textcolor{OliveGreen}{yes} &  \\
C29 & structure: made of heart\_\allowbreak{}rate\_\allowbreak{}nonanginal\_\allowbreak{}gate,\allowbreak{} restecg2\_\allowbreak{}relay and inputs age, sex, cp, trestbps, chol, fbs, thalach, oldpeak, slope, ca, thal; added &  & log-odds n/a, OR n/a & 242 & train & \textcolor{OliveGreen}{yes} &  \\
C30 & direction: chest pain type: chest pain type other than non-anginal chest pain to non-anginal chest pain (through this part) \newline\textit{shape: categorical; highest to lowest: chest pain type other than non-anginal chest pain, non-anginal chest pain} & 32\% to 16\% & log-odds -0.79, OR 0.46; range across patients -1.14 to -0.21 & 172, 70 & train; test unclear (AUC gain 0.008 [-0.012, 0.030], n=61) & \textcolor{OliveGreen}{yes} & on the training data only \\
C31 & direction: maximum heart rate: 118 beats per minute to 175 beats per minute (through this part) \newline\textit{shape: hinge at 168 beats per minute; 70\% of patients same; mixed (but steady for patients with non-anginal chest pain); one point with C36 (median fold 168 beats per minute, typical patient 166 beats per minute, 10th-90th percentile 168 beats per minute to 168 beats per minute)} & 32\% to 29\% & log-odds -0.34, OR 0.71; range across patients -0.93 to -0.12 & 242 & train; test unclear (AUC gain 0.008 [0.000, 0.022], n=61) & \textcolor{OliveGreen}{yes} & slightly; on the training data only \\
C32 & direction: thallium stress test result: thallium stress test result other than a fixed defect on the thallium scan to a fixed defect on the thallium scan (through this part) \newline\textit{shape: categorical; highest to lowest: a fixed defect on the thallium scan, thallium stress test result other than a fixed defect on the thallium scan} & 2\% to 4\% (typical patient this part acts on) & log-odds 0.24, OR 1.27; range across patients 0.00 to 0.82 & 227, 15 (small n) & train; test unclear (AUC gain 0.001 [0.000, 0.006], n=61) & \textcolor{OliveGreen}{yes} & slightly; the model has little data on this; on the training data only \\
C33 & interaction: effect of chest pain type (chest pain type other than non-anginal chest pain to non-anginal chest pain) by thallium stress test result: a fixed defect on the thallium scan minus thallium stress test result other than a fixed defect on the thallium scan; stronger in thallium stress test result other than a fixed defect on the thallium scan \newline\textit{shape: categorical; highest to lowest: chest pain type other than non-anginal chest pain, non-anginal chest pain} & a fixed defect on the thallium scan: 12\% to 2\%; thallium stress test result other than a fixed defect on the thallium scan: 32\% to 3\% & log-odds difference 0.79, OR 2.20 & 15, 227 (small n) & train; test agrees (AUC gain 0.001 [0.000, 0.006], n=61); group difference on test 0.73 [0.60, 0.85] & \textcolor{OliveGreen}{yes} & the model has little data on this \\
C34 & interaction: effect of maximum heart rate (118 beats per minute to 175 beats per minute) by thallium stress test result: a fixed defect on the thallium scan minus thallium stress test result other than a fixed defect on the thallium scan; stronger in thallium stress test result other than a fixed defect on the thallium scan \newline\textit{shape: accelerating; 70\% of patients same; mixed (but steady for patients with non-anginal chest pain)} & a fixed defect on the thallium scan: 14\% to 10\%; thallium stress test result other than a fixed defect on the thallium scan: 37\% to 26\% & log-odds difference 0.30, OR 1.35 & 15, 227 (small n) & train; test agrees (AUC gain 0.001 [0.000, 0.006], n=61); group difference on test 0.34 [0.26, 0.44] & \textcolor{OliveGreen}{yes} & the model has little data on this \\
C35 & interaction: effect of maximum heart rate (118 beats per minute to 175 beats per minute) by chest pain type: non-anginal chest pain minus chest pain type other than non-anginal chest pain; stronger in non-anginal chest pain \newline\textit{shape: steady; 97\% of patients same; verified; risk saturates near 0\%} & non-anginal chest pain: 7\% to 2\%; chest pain type other than non-anginal chest pain: 37\% to 26\% & log-odds difference -0.80, OR 0.45 & 70, 172 & train; test agrees (AUC gain 0.008 [-0.012, 0.030], n=61); group difference on test -0.82 [-0.83, -0.81] & \textcolor{OliveGreen}{yes} & clearly \\
C36 & threshold: maximum heart rate: fold at about 168 beats per minute (starts) for patients with no chest pain symptoms (asymptomatic) \newline\textit{shape: hinge at 168 beats per minute; 100\% of patients same; verified; risk saturates near 100\%; one point with C36 (median fold 168 beats per minute, typical patient 166 beats per minute, 10th-90th percentile 168 beats per minute to 168 beats per minute)} & 99\% to 99\% (typical patient with a fold (patients with no chest pain symptoms (asymptomatic))) & fold 168 beats per minute (IQR 168 to 168); past it log-odds -0.23, OR 0.79; range across patients -0.23 to -0.23 & 99, 14 (small n) & train; test unclear (AUC gain 0.008 [0.000, 0.022], n=61) & \textcolor{OliveGreen}{yes} & slightly; the model has little data on this; on the training data only \\
\bottomrule
\end{tabularx}}
\par{\tiny\noindent
\begin{tabularx}{\linewidth}{@{}l>{\raggedright\arraybackslash}X>{\raggedright\arraybackslash}p{0.11\linewidth}>{\raggedright\arraybackslash}p{0.16\linewidth}>{\raggedright\arraybackslash}p{0.05\linewidth}>{\raggedright\arraybackslash}p{0.15\linewidth}l>{\raggedright\arraybackslash}p{0.08\linewidth}@{}}
\toprule
Claim & What is checked & Risk, typical patient & Estimate (log-odds; OR); range across patients & n & Test check & Verified & Words \\
\midrule
C37 & subgroup: who this part acts on: 104 of 242 patients & 5\% to 2\% (typical patient this part acts on (risk without the part, then with it)) & share 0.43; log-odds -0.67, OR 0.51; range across patients -1.22 to -0.05 & 104, 138 & train; test unclear (AUC gain 0.018 [0.000, 0.044], n=61) & \textcolor{OliveGreen}{yes} & slightly; on the training data only \\
C38 & structure: its formula reproduces it on 242 training patients (largest difference 4.44089e-16) &  & log-odds 4e-16, OR n/a & 242 & train & \textcolor{OliveGreen}{yes} &  \\
C39 & direction: resting ECG: resting ECG other than signs of left ventricular hypertrophy on the resting ECG to signs of left ventricular hypertrophy on the resting ECG (through this part) \newline\textit{shape: categorical; highest to lowest: resting ECG other than signs of left ventricular hypertrophy on the resting ECG, signs of left ventricular hypertrophy on the resting ECG} & 32\% to 31\% & log-odds -0.02, OR 0.98; range across patients -0.02 to -0.02 & 123, 119 & train; test unclear (AUC gain 0.000 [0.000, 0.000], n=61) & \textcolor{OliveGreen}{yes} & slightly; on the training data only \\
C40 & subgroup: who this part acts on: 119 of 242 patients & 85\% to 84\% (typical patient this part acts on (risk without the part, then with it)) & share 0.49; log-odds -0.02, OR 0.98; range across patients -0.02 to -0.02 & 119, 123 & train; test unclear (AUC gain 0.000 [0.000, 0.000], n=61) & \textcolor{BrickRed}{no} & slightly; on the training data only \\
C41 & null: this part barely changes predictions (switched off) &  & max change 0.005 (max $|$change in predicted risk$|$) & 242 & train & \textcolor{OliveGreen}{yes} &  \\
C42 & structure: its formula reproduces it on 242 training patients (largest difference 0.0) &  & log-odds 0.00, OR n/a & 242 & train & \textcolor{OliveGreen}{yes} &  \\
\bottomrule
\end{tabularx}}
\paragraph{Not supported by the model.}
\begin{itemize}\setlength\itemsep{0pt}\small
\item C40: subgroup claim about resting ECG in part ann\_\allowbreak{}29: not supported by the model (its effect is negligible).
\end{itemize}

\subsection{Coverage}\label{app:heart-s2-c2:coverage}
Coverage is computed by the current code (the function behind the \texttt{get\_coverage} tool) under these conventions. \emph{Candidate nodes} are all nodes of the explained model state except the output nodes. A candidate node is \emph{covered} when some annotation of the hierarchy (a leaf, or a composition together with its descendants) contains it and all its enabled outgoing edges. \emph{Exact Null} parts, which have no effect on any output for any input, count as covered: nodes with no path to an output over enabled, non-zero-weight edges (disconnected inputs and dead ends), and edges that have zero weight or touch such a node. Compositional coverage is the fraction of required compositions present: every annotation except a single global root containing every output must have a parent; a parentless, childless annotation made only of Null nodes is exempt.

\begin{itemize}
\item \textbf{Null, disconnected input (8):} \texttt{-10} (fbs:0), \texttt{-12} (restecg:0), \texttt{-13} (restecg:1), \texttt{-16} (exang:0), \texttt{-17} (exang:1), \texttt{-21} (slope:3), \texttt{-23} (thal:3), \texttt{-24} (thal:4).
\item \textbf{Uncovered nodes:} none.
\item \textbf{Uncovered edges:} none.
\item \textbf{Missing compositions:} none.
\end{itemize}

\subsection{Annotation hierarchy}\label{app:heart-s2-c2:hierarchy}
The hierarchy has 1 root; the global root is heart\_\allowbreak{}disease\_\allowbreak{}logit. Tags (A1, A2, \ldots) are used in the figures.
\begingroup\parindent=0pt\parskip=1pt
\hspace*{0.0em}$\triangleright$~A1: heart\_\allowbreak{}disease\_\allowbreak{}logit (composition, \cref{app:heart-s2-c2:ann-heart-disease-logit})\par
\hspace*{1.5em}$\triangleright$~A2: heart\_\allowbreak{}rate\_\allowbreak{}nonanginal\_\allowbreak{}gate (leaf, \cref{app:heart-s2-c2:ann-heart-rate-nonanginal-gate})\par
\hspace*{1.5em}$\triangleright$~A3: restecg2\_\allowbreak{}relay (leaf, \cref{app:heart-s2-c2:ann-restecg2-relay})\par
\endgroup

\subsection{Network figures}\label{app:heart-s2-c2:figures}
\begin{figure}[htbp]
\centering
\input{\explaneatcertdir{heart-s2-c2}heart-s2-c2-fig-full.tex}
\caption{heart-s2-c2: full view of the explained model state.}
\label{app:heart-s2-c2:fig-full}
\end{figure}
\begin{figure}[htbp]
\centering
\input{\explaneatcertdir{heart-s2-c2}heart-s2-c2-fig-level-0.tex}
\caption{heart-s2-c2: collapsed view at hierarchy level 0 (fully recomposed).}
\label{app:heart-s2-c2:fig-level-0}
\end{figure}
\begin{figure}[htbp]
\centering
\input{\explaneatcertdir{heart-s2-c2}heart-s2-c2-fig-level-1.tex}
\caption{heart-s2-c2: collapsed view at hierarchy level 1.}
\label{app:heart-s2-c2:fig-level-1}
\end{figure}

\subsection{Refactoring operations}\label{app:heart-s2-c2:operations}
The explained model state is the phenotype after these operations, replayed in order. All are function-preserving; the largest change in the model output over the 242 training rows is $2.22\times 10^{-16}$ (tolerance $1.00\times 10^{-9}$).
\begin{itemize}
\item \#26: split \texttt{-15} into \texttt{-15\_\allowbreak{}a}, \texttt{-15\_\allowbreak{}b}.
\item \#27: split \texttt{-25} into \texttt{-25\_\allowbreak{}a}, \texttt{-25\_\allowbreak{}b}.
\end{itemize}

\subsection{Annotations}\label{app:heart-s2-c2:annotations}
\subsubsection{\texorpdfstring{A1: heart\_\allowbreak{}disease\_\allowbreak{}logit}{A1: heart disease logit}}\label{app:heart-s2-c2:ann-heart-disease-logit}
% annotation id ann_30, operation seq 30
\textbf{Kind:} Composition annotation (depth 0); children A2 (\cref{app:heart-s2-c2:ann-heart-rate-nonanginal-gate}), A3 (\cref{app:heart-s2-c2:ann-restecg2-relay}).

\begin{itemize}
\item \textbf{Entry nodes:} \texttt{-1} (age), \texttt{-2} (sex:0), \texttt{-3} (sex:1), \texttt{-4} (cp:1), \texttt{-5} (cp:2), \texttt{-7} (cp:4), \texttt{-8} (trestbps), \texttt{-9} (chol), \texttt{-10} (fbs:0), \texttt{-11} (fbs:1), \texttt{-12} (restecg:0), \texttt{-13} (restecg:1), \texttt{-15\_\allowbreak{}a} (thalach\_\allowbreak{}a), \texttt{-16} (exang:0), \texttt{-17} (exang:1), \texttt{-18} (oldpeak), \texttt{-19} (slope:1), \texttt{-20} (slope:2), \texttt{-21} (slope:3), \texttt{-22} (ca), \texttt{-23} (thal:3), \texttt{-24} (thal:4), \texttt{-25\_\allowbreak{}a} (thal:6\_\allowbreak{}a), \texttt{-26} (thal:7).
\item \textbf{Exit nodes:} \texttt{0}.
\item \textbf{Internal nodes:} none.
\item \textbf{Edges (weight):} \texttt{-1}$\to$\texttt{0}~(0.0279), \texttt{-2}$\to$\texttt{0}~(-0.771), \texttt{-3}$\to$\texttt{0}~(0.34), \texttt{-4}$\to$\texttt{0}~(-0.369), \texttt{-5}$\to$\texttt{0}~(0.264), \texttt{-7}$\to$\texttt{0}~(0.869), \texttt{-8}$\to$\texttt{0}~(0.492), \texttt{-9}$\to$\texttt{0}~(0.535), \texttt{-11}$\to$\texttt{0}~(-0.222), \texttt{-18}$\to$\texttt{0}~(0.556), \texttt{-19}$\to$\texttt{0}~(-0.469), \texttt{-20}$\to$\texttt{0}~(0.634), \texttt{-22}$\to$\texttt{0}~(1.78), \texttt{-26}$\to$\texttt{0}~(0.501), \texttt{3607}$\to$\texttt{0}~(-0.0209), \texttt{3842}$\to$\texttt{0}~(-0.519), \texttt{-15\_\allowbreak{}a}$\to$\texttt{0}~(-0.153), \texttt{-25\_\allowbreak{}a}$\to$\texttt{0}~(-0.3).
\item \textbf{With its descendants:} 33 nodes.
\item \textbf{Covers} (coverage, own and as part of the hierarchy): 32 candidate nodes.
\end{itemize}

\textbf{Hypothesis.} Root: heart-disease logit at output 0 (sigmoid). Essentially logistic regression: 16 direct edges (ca +1.78, cp:4 +0.87, sex:0 -0.77, slope:2 +0.63, oldpeak +0.56, chol +0.54, thal:7 +0.50, trestbps +0.49, slope:1 -0.47, cp:1 -0.37, sex:1 +0.34, thal:6 -0.30, cp:2 +0.26, fbs:1 -0.22, thalach -0.15, age +0.03) carry logit SD 3.66 and train AUC 0.9343 alone (L2 logistic regression on all 26 inputs: 0.9355). The two children are added to the logit: heart\_\allowbreak{}rate\_\allowbreak{}nonanginal\_\allowbreak{}gate (x -0.519; +0.006 AUC, the only consequential nonlinearity) and restecg2\_\allowbreak{}relay (x -0.021, negligible). Eight inputs are unconnected, zero effect: thal:3, thal:4, slope:3, exang:0, exang:1, restecg:0, restecg:1, fbs:0 (exercise angina is ignored). Evidence (train, n=242): composed closed form; sensitivity and mean $|$SHAP$|$ both rank ca first (0.16), then cp:4, sex:0, slope:2; full-model train AUC 0.9403, accuracy 0.897.

\textbf{Closed form} (children as named functions):
{\small\begingroup\allowdisplaybreaks
\begin{align*}
y_{0} &= \sigma\bigl(z_{0}\bigr) \\
z_{0} &= 0.324 + 0.501 x_{thal:7} + 0.556 x_{oldpeak} + 0.34 x_{sex:1} + 0.0279 x_{age} \\
&\qquad + 0.634 x_{slope:2} + 0.264 x_{cp:2} + 0.535 x_{chol} + 0.869 x_{cp:4} + 1.78 x_{ca} \\
&\qquad + 0.492 x_{trestbps} - 0.771 x_{sex:0} - 0.153 x_{thalach\_a} \\
&\qquad - 0.0209 \operatorname{restecg2\_relay}\bigl(x_{restecg:2}\bigr) - 0.369 x_{cp:1} \\
&\qquad - 0.519 \operatorname{heart\_rate\_nonanginal\_gate}\bigl(x_{cp:3}, x_{thalach\_b}, \\
&\qquad x_{thal:6\_b}\bigr) - 0.3 x_{thal:6\_a} - 0.222 x_{fbs:1} - 0.469 x_{slope:1}
\end{align*}\endgroup}

Expanded over the children:
{\small\begingroup\allowdisplaybreaks
\begin{align*}
y_{0} &= \sigma\bigl(z_{0}\bigr) \\
z_{0} &= 0.324 + 0.501 x_{thal:7} + 0.556 x_{oldpeak} + 0.34 x_{sex:1} + 0.0279 x_{age} \\
&\qquad + 0.634 x_{slope:2} + 0.264 x_{cp:2} + 0.535 x_{chol} + 0.869 x_{cp:4} + 1.78 x_{ca} \\
&\qquad + 0.492 x_{trestbps} - 0.771 x_{sex:0} - 0.153 x_{thalach\_a} \\
&\qquad - 0.0209 \operatorname{relu}\bigl(-0.016 + 0.997 x_{restecg:2}\bigr) - 0.369 x_{cp:1} \\
&\qquad - 0.519 \operatorname{relu}\bigl(-0.295 + 1.02 \operatorname{relu}\bigl(-0.347 \\
&\qquad + 0.688 x_{thalach\_b} + 1.03 \operatorname{relu}\bigl(0.571 + 0.968 x_{cp:3}\bigr) \\
&\qquad - 0.374 x_{thal:6\_b}\bigr)\bigr) - 0.3 x_{thal:6\_a} - 0.222 x_{fbs:1} - 0.469 x_{slope:1}
\end{align*}\endgroup}

where $y_{0}$ is node \texttt{0}; $z_{i}$, where shown, is the pre-activation of a sigmoid output $y_{i}$; weights and biases are shown to 3 significant figures.
 On the training split (n = 242), the displayed closed form differs from the forward pass by at most $8.94\times 10^{-4}$, and the exact function by at most $2.22\times 10^{-16}$.

\textbf{Evidence} (2 items; full payloads in the JSON record):
\begin{itemize}
\item \textsc{analytical}, \texttt{analytical\_\allowbreak{}formula}; data split: n/a (analytical). Root closed form (get\_\allowbreak{}formula force=True): direct logistic terms plus the two child functions
\item \textsc{empirical}, \texttt{ablation\_\allowbreak{}study}; data split: train. Train split, full sample; ablations and linear replacements on in-memory copies (no graph edits); SHAP on all 242 train rows
\end{itemize}

\textbf{Authorship:} drafted by AI-assisted draft on 2026-11-10; verified by nobody (unverified).

\subsubsection{\texorpdfstring{A2: heart\_\allowbreak{}rate\_\allowbreak{}nonanginal\_\allowbreak{}gate}{A2: heart rate nonanginal gate}}\label{app:heart-s2-c2:ann-heart-rate-nonanginal-gate}
% annotation id ann_28, operation seq 28
\textbf{Kind:} Leaf annotation (depth 1); parent A1 (\cref{app:heart-s2-c2:ann-heart-disease-logit}).

\begin{itemize}
\item \textbf{Entry nodes:} \texttt{-6} (cp:3), \texttt{-15\_\allowbreak{}b} (thalach\_\allowbreak{}b), \texttt{-25\_\allowbreak{}b} (thal:6\_\allowbreak{}b).
\item \textbf{Exit nodes:} \texttt{3842}.
\item \textbf{Internal nodes:} \texttt{2942}, \texttt{2472}.
\item \textbf{Edges (weight):} \texttt{-6}$\to$\texttt{2942}~(0.968), \texttt{-15\_\allowbreak{}b}$\to$\texttt{2472}~(0.688), \texttt{-25\_\allowbreak{}b}$\to$\texttt{2472}~(-0.374), \texttt{2942}$\to$\texttt{2472}~(1.03), \texttt{2472}$\to$\texttt{3842}~(1.02).
\item \textbf{Covers} (coverage, own and as part of the hierarchy): 5 candidate nodes.
\end{itemize}

\textbf{Hypothesis.} Protective heart-rate gate, gated by chest-pain type and thallium. 2942 = ReLU(0.968 cp:3 + 0.571) is a non-anginal-pain indicator (2.09 for the 70 cp:3 patients, 0 otherwise). 3842 = ReLU(1.021*\allowbreak{}ReLU(0.688 thalach - 0.374 thal:6 + 1.028*2942 - 0.347) - 0.295) is read out -0.519 into the logit. The gate opens once max heart rate passes a threshold set by the other two inputs: about -2.3 SD with non-anginal pain (open for 66 of 68), about +0.78 SD for other pain types without a fixed defect (38 of 159), and never with a fixed thallium defect (0 of 15). Once open, each SD of thalach lowers the logit by 0.36 more (total -0.52 vs -0.15 closed), and non-anginal pain lowers it by about 1.1. Evidence (train, n=242): closed form from get\_\allowbreak{}formula; gate open for 104 patients (disease rate 23\% vs 60\% when closed); logit contribution SD 0.45; removing it drops train AUC 0.9403 -$>$ 0.9344, a linear refit on cp:3, thalach, thal:6 recovers 0.9378 (R2 0.87), so most of its value is the gating rather than the linear trend. It is the model's only consequential nonlinearity.

\textbf{Closed form:}
{\small\begingroup\allowdisplaybreaks
\begin{align*}
y_{0} &= \operatorname{relu}\bigl(-0.295 + 1.02 \operatorname{relu}\bigl(-0.347 \\
&\qquad + 0.688 x_{thalach\_b} + 1.03 \operatorname{relu}\bigl(0.571 + 0.968 x_{cp:3}\bigr) \\
&\qquad - 0.374 x_{thal:6\_b}\bigr)\bigr)
\end{align*}\endgroup}

where $y_{0}$ is node \texttt{3842}; $z_{i}$, where shown, is the pre-activation of a sigmoid output $y_{i}$; weights and biases are shown to 3 significant figures.
 On the training split (n = 242), the displayed closed form differs from the forward pass by at most 0.00534, and the exact function by at most $4.44\times 10^{-16}$.

\textbf{Evidence} (2 items; full payloads in the JSON record):
\begin{itemize}
\item \textsc{analytical}, \texttt{analytical\_\allowbreak{}formula}; data split: n/a (analytical). Closed form from get\_\allowbreak{}formula (force=True)
\item \textsc{empirical}, \texttt{empirical\_\allowbreak{}viz} (heatmap); data split: train. Train split, full sample (242); local forward passes on in-memory copies for ablations
\end{itemize}

\textbf{Authorship:} drafted by AI-assisted draft on 2026-11-10; verified by nobody (unverified).

\subsubsection{\texorpdfstring{A3: restecg2\_\allowbreak{}relay}{A3: restecg2 relay}}\label{app:heart-s2-c2:ann-restecg2-relay}
% annotation id ann_29, operation seq 29
\textbf{Kind:} Leaf annotation (depth 1); parent A1 (\cref{app:heart-s2-c2:ann-heart-disease-logit}).

\begin{itemize}
\item \textbf{Entry nodes:} \texttt{-14} (restecg:2).
\item \textbf{Exit nodes:} \texttt{3607}.
\item \textbf{Internal nodes:} none.
\item \textbf{Edges (weight):} \texttt{-14}$\to$\texttt{3607}~(0.997).
\item \textbf{Covers} (coverage, own and as part of the hierarchy): 1 candidate node.
\end{itemize}

\textbf{Hypothesis.} restecg:2 relay: 3607 = ReLU(0.997 restecg:2 - 0.016) is 0.998 for the 119 patients with left ventricular hypertrophy on the resting ECG and 0 for the other 123, i.e. exactly affine in the binary input on the data. Read out -0.021, it lowers the logit by 0.02 for LVH: negligible. Evidence (train, n=242): closed form; entry-exit correlation 1.0; logit SD 0.01; removing it changes train AUC 0.9403 -$>$ 0.9401.

\textbf{Closed form:}
{\small\begingroup\allowdisplaybreaks
\begin{align*}
y_{0} &= \operatorname{relu}\bigl(-0.016 + 0.997 x_{restecg:2}\bigr)
\end{align*}\endgroup}

where $y_{0}$ is node \texttt{3607}; $z_{i}$, where shown, is the pre-activation of a sigmoid output $y_{i}$; weights and biases are shown to 3 significant figures.
 On the training split (n = 242), the displayed closed form differs from the forward pass by at most $1.39\times 10^{-4}$, and the exact function by at most 0.

\textbf{Evidence} (2 items; full payloads in the JSON record):
\begin{itemize}
\item \textsc{analytical}, \texttt{analytical\_\allowbreak{}formula}; data split: n/a (analytical). Closed form from get\_\allowbreak{}formula (force=True)
\item \textsc{empirical}, \texttt{empirical\_\allowbreak{}viz} (line\_\allowbreak{}plot); data split: train. Train split, full sample (242); local forward passes on in-memory copies for ablations
\end{itemize}

\textbf{Authorship:} drafted by AI-assisted draft on 2026-11-10; verified by nobody (unverified).

\subsection{Audit against the definitions}\label{app:heart-s2-c2:audit}
Automatic findings on the stored explanation (codes as in the JSON record's \texttt{issues}); they flag what an examiner should check, not verdicts.
\begin{itemize}
\item \textsc{warning} \texttt{unverified}: No human verification is recorded for this explanation.
\end{itemize}
\end{drafted}
